\documentclass{exam-zh}
\usepackage{siunitx}
\usepackage{array} % 用于指定列宽

% 定义新的列类型,指定列宽为 0.80cm 且内容居中
\newcolumntype{C}{>{\centering\arraybackslash}p{0.80cm}} %设置表格宽度

% 显示答案
\examsetup{
    % title/title-format={\huge},
    page/size=a4paper,
    paren/show-paren=true,
    paren/show-answer=true,
    paren/text-color=red,
    fillin/show-answer=true,
    fillin/text-color=blue,
    solution/text-color=blue,
    solution/show-solution=show-stay,
    page/foot-content=第;页，共;页,
}


\ExamPrintAnswerSet {
page/size=a3paper,
solution/show-solution=show-stay,
paren/show-paren=true,
}

% 不显示答案（如需隐藏答案,可注释上方“显示答案”配置,启用下方代码）]


% \examsetup{
%   page/size=a4paper,             % 试卷页面尺寸设为A4
%   paren/show-paren=true,        % 显示选择题答案括号（如"( )"）
%   paren/show-answer=false,      % 隐藏选择题括号内答案
%   paren/text-color=red,         % 选择题答案文字颜色（若显示则为红色）
%   fillin/show-answer=false,     % 隐藏填空题答案
%   fillin/text-color=blue,       % 填空题答案文字颜色（若显示则为蓝色）
%   fillin/no-answer-type=none,   % 填空题隐藏答案时无占位符（去除黑三角）
%   solution/show-solution=hide,  % 隐藏解答题解答内容
%   % solution/blank-type=manual,  % （可选）隐藏解答时手动加空白
%   % solution/blank-vsep=60ex plus 1ex minus 1ex % （可选）控制空白高度
% }
% % \ExamPrintAnswerSet[\geometry{showframe}]{
% % page/size=a3paper,
% % solution/show-solution=show-stay,
% % paren/show-paren=true,
% % }

\everymath{\displaystyle}

\usepackage{tasks}
\usepackage{tabularx}

\begin{document} 

\title{杭州市中职数学思维等级证书\\红段测试题}
\maketitle

\information{
考试时间: 45分钟,
满分: 100分
}

\section{选择题(每题4分,共60分)}

\begin{question}
     已知集合 $A = \{x \in \mathbf{N} \mid x < 5\}$,集合 $B = \{3,4,6\}$,则 $A \cup B =$ \paren[A]
    \begin{choices}
        \item $\{3,4,6\}$ 
        \item $\{0,1,2,3,4,6\}$ 
        \item $\{1,2,3,4,5,6\}$ 
        \item $\{0,1,2,3,4,5,6\}$
    \end{choices}
\end{question}

\begin{question}
    已知 $a > b$,$c > 0$,下列结论正确的是 \paren[B]
    \begin{choices}
        \item $a - c < b - c$
        \item $ac > bc$
        \item $a + c < b + c$
        \item $\dfrac{a}{c} < \dfrac{b}{c}$
    \end{choices}
\end{question}
\begin{question}
    倾斜角为 $\dfrac{\pi}{6}$,在 $y$ 轴上的截距为 $4$,求直线的方程 \paren[A]
    \begin{choices}
        \item $x - \sqrt{3}y + 4\sqrt{3} = 0$
        \item $x + \sqrt{3}y + 4\sqrt{3} = 0$
        \item $x - \sqrt{3}y - 4\sqrt{3} = 0$
        \item $x + \sqrt{3}y - 4\sqrt{3} = 0$
    \end{choices}
\end{question}

\begin{question}
    已知 $x \in (0,\pi)$,则 $\cos x = -\dfrac{1}{2}$ 的解为 \paren[B]
    \begin{choices}
        \item $\dfrac{\pi}{3}$
        \item $\dfrac{2\pi}{3}$
        \item $\dfrac{\pi}{6}$
        \item $\dfrac{5\pi}{6}$
    \end{choices}
\end{question}

\begin{question}
    $\tan x > 0$ 是 $x$ 在第一或第三象限的 \paren[C]
    \begin{choices}
        \item 充分不必要条件
        \item 必要不充分条件
        \item 充要条件
        \item 既不充分又不必要条件
    \end{choices}
\end{question}

\begin{question}
    函数 $f(x)=\sqrt{\lg(x + 3)}+\dfrac{1}{\sqrt{3 - x}}$ 的定义域为\paren[C]
    \begin{choices}
        \item $(-3,3)$
        \item $(-2,3]$
        \item $[-2,3)$
        \item $(-3,4]$
    \end{choices}
\end{question}

\begin{question}
    方程 $|\sqrt{(x - 3)^2 + y^2}-\sqrt{(x + 3)^2 + y^2}|=4$ 表示的\paren[D]
    \begin{choices}
        \item 线段
        \item 椭圆
        \item 直线
        \item 双曲线
    \end{choices}
\end{question}




\begin{question}
    二次函数 $y = x^2 - 2x - 4$ 在区间 $[-1,2]$ 上的最大值、最小值分别为\paren[A]
    \begin{choices}
        \item $-1, -5$
        \item $-1, -4$
        \item $-4, -5$
        \item $-5, -4$
    \end{choices}
    \end{question}
    
    \begin{question}
    如果角 $\alpha$ 的终边过点 $(5,m)$,且 $\cos\alpha=\dfrac{5}{13}$,则 $\tan(\pi - \alpha)=$\paren[A]
    \begin{choices}
        \item $-\dfrac{12}{5}$
        \item $\pm\dfrac{12}{5}$
        \item $\dfrac{5}{12}$
        \item $\pm\dfrac{5}{12}$
    \end{choices}
    \end{question}
    
    \begin{question}
    已知向量 $\vec{a}=(4,-2)$,$\vec{b}=(2,3)$,则 $|\vec{a} + 2\vec{b}|=$\paren[B]
    \begin{choices}
        \item $(8,4)$
        \item $4\sqrt{5}$
        \item $(8,8)$
        \item $8\sqrt{2}$
    \end{choices}
    \end{question}


\end{document}
% xelatex -shell-escape -interaction=batchmode 
